\subsection{Noether normalization, field extensions and the graded vertex}\label{algebra-proof-026-noether-normalization-field-extensions-and-the-graded-vertex}

Authority: \texttt{algebra.tex} 27920--28447 at commit \texttt{a044\allowbreak{}46e5\allowbreak{}7ec1\allowbreak{}fbc2\allowbreak{}52a8\allowbreak{}71af\allowbreak{}cec7\allowbreak{}752f\allowbreak{}b280\allowbreak{}7b14}, SHA-256 \texttt{FA8B\allowbreak{}B92E\allowbreak{}58A4\allowbreak{}F78A\allowbreak{}2BD0\allowbreak{}1B3B\allowbreak{}6A4A\allowbreak{}87DE\allowbreak{}0A0D\allowbreak{}279F\allowbreak{}5DD9\allowbreak{}0641\allowbreak{}B574\allowbreak{}DD5F\allowbreak{}BFFF\allowbreak{}A4F3}. The complete current interval equals this authority byte for byte. There are no prior canonical operations or added environments in the interval.

These complete arguments and further consequences are separate editorial evidence. They preserve the original polynomials, automorphisms, units, prime labels and torsion. Neither translation branch is replaced by these expanded proofs. Proved consequences carry no novelty claim.

\subsubsection{Finite supports and explicit positive weights}\label{algebra-proof-026-finite-supports-and-explicit-positive-weights}

Source 27927--27960. Let \(N\subset\mathbb Z_{\geq0}^n\) be the original finite nonempty support, \(n\geq1\), and retain \(A_i=\max_{\nu\in N}\nu_i-\min_{\nu\in N}\nu_i\). The required inequalities are indexed by \(i=2,\ldots,n\); there is no \(e_0\). This resolves OCC-12271 and OCC-00545.

An explicit choice is \[
 e_n=1,\qquad
 e_i=1+\sum_{j=i+1}^n A_je_j
     =\prod_{j=i+1}^n(A_j+1)\quad(1\leq i<n).
\] The second equality follows downwards from \(\sum_{j=i+1}^n A_j\prod_{\ell=j+1}^n(A_\ell+1)
=\prod_{\ell=i+1}^n(A_\ell+1)-1\), a telescoping sum. Every weight is a positive integer. If \(\nu,\nu'\) first differ in coordinate \(i\), then \[
 |e_i(\nu_i-\nu_i')|\geq e_i
 >\sum_{j>i}A_je_j
 \geq\left|\sum_{j>i}e_j(\nu_j-\nu_j')\right|.
\] The first differing coordinate therefore determines the sign of the weighted difference, which cannot vanish. This proves injectivity on the entire original support, retaining its actual coordinate ranges. For \(n=1\), \(e_1=1\) gives the assertion directly, with no tail inequality. The optional phrase ``with \(a_\nu\in R\) not zero'' is mathematically and grammatically intelligible; OCC-00546 does not warrant a required source edit.

For any finite family of nonconstant polynomials in these same variables, choose the weights using the union of their finite supports. This one choice separates the support of every member simultaneously.

\subsubsection{The original substitutions, full coefficients and one relation}\label{algebra-proof-026-the-original-substitutions-full-coefficients-and-one-relation}

Source 27962--28022. Write the original polynomial \(g=\sum_{\nu\in N}a_\nu x^\nu\), with all listed \(a_\nu\ne0\). Keep the automorphism \[
 \sigma(x_i)=x_i+x_n^{e_i}\ (i<n),\qquad \sigma(x_n)=x_n.
\] Its inverse sends \(x_i\) to \(x_i-x_n^{e_i}\) for \(i<n\) and fixes \(x_n\); the two composites fix each generator. The full expansion of its \(\nu\)-summand is \[
 a_\nu\sum_{\substack{0\leq b_i\leq\nu_i\\i<n}}
 \left(\prod_{i<n}\binom{\nu_i}{b_i}x_i^{b_i}\right)
 x_n^{\,\nu_n+\sum_{i<n}e_i(\nu_i-b_i)}.
\] The term with every \(b_i=0\) is \(a_\nu x_n^{e\cdot\nu}\). Every other term has strictly smaller \(x_n\)-degree, since all weights are positive. The unique support vector \(\nu^*\) maximizing \(e\cdot\nu\) therefore gives \[
 \sigma(g)=a_{\nu^*}x_n^{d}
              +\sum_{j=0}^{d-1}c_j(x_1,\ldots,x_{n-1})x_n^j,
 \qquad d=e\cdot\nu^*>0 .
\] This is valid over arbitrary coefficient rings, including rings with nilpotents, because the leading coefficient is the single original nonzero coefficient \(a_{\nu^*}\); no cancellation or division has been used. All other binomial summands remain in the displayed expansion.

For a ring map \(R\to R'\), the same integer substitution and its inverse commute with coefficient base change. Some original coefficients may vanish. If the image polynomial is nonconstant, its surviving support is a subset of \(N\), so the same weights still separate it and its new leading coefficient is the image of one specified surviving original coefficient. If all nonconstant terms vanish, the image is constant or zero, and no positive-degree conclusion is asserted. This also proves the simultaneous assertion for the finite family above under any base change, with its precise surviving-support boundary.

Over the original field \(k\), choose nonzero \(f\in I\) with \(I\) proper. Then \(f\) is nonconstant, since a nonzero constant is a unit. Put \(y_i=x_i-x_n^{e_i}\) for \(i<n\). These actual polynomials have integer coefficients, and \(k[y_1,\ldots,y_{n-1},x_n]=k[x_1,\ldots,x_n]\), with inverse expressions \(x_i=y_i+x_n^{e_i}\). In the quotient by \(f\), retain the complete equation \[
 a x_n^d+\sum_{j<d}c_j(y_1,\ldots,y_{n-1})x_n^j=0,\qquad a\in k^\times.
\] It gives the explicit recurrence \(x_n^d=-a^{-1}\sum_{j<d}c_j(y)x_n^j\). Induction on the exponent reduces every power to the span of \(1,x_n,\ldots,x_n^{d-1}\). Hence the quotient by \((f)\), and then its quotient \(S\), is finite over the actual image of \(k[y_1,\ldots,y_{n-1}]\). The inverse unit \(a^{-1}\) has not been dropped from the calculation. For \(n=1\) there are no \(y_i\), and this is the same univariate proof. For \(n=0\) a field has no nonzero proper ideal, so the premise does not occur. OCC-00547 is the subject--verb correction ``there exist integers''.

\subsubsection{The full induction and a neighborhood of the original point}\label{algebra-proof-026-the-full-induction-and-a-neighborhood-of-the-original-point}

Source 28024--28071. For \(n=0\), the proper-ideal quotient is \(k\); take \(r=0\). For \(I=0\), take the original variables and \(r=n\). Otherwise use the preceding \(n-1\) integer polynomials and let \(S'\) be their actual image subalgebra of \(S\). This is nonzero and is generated by \(n-1\) elements. Apply the induction hypothesis to a polynomial presentation of that image. Its generators \(z_j\) are integer polynomials in the \(y_i\); substituting their actual integer-polynomial expressions in the original \(x_i\) retains integer coefficients. The composition \(k[z_1,\ldots,z_r]\to S'\to S\) is injective and finite. To check the latter directly, products of a finite generating family for \(S'\) over \(k[z]\) and a finite generating family for \(S\) over \(S'\) generate \(S\) over \(k[z]\).

The dimension statement follows from injective integral dimension equality and \(\dim k[z_1,\ldots,z_r]=r\). The integer-polynomial choice follows from the explicit construction, not from that dimension lemma. Thus OCC-12272 and OCC-00548 replace the misleading ``last assertion'' by the assertion that \(r\) is the dimension of \(S\).

For a point \(x\leftrightarrow\mathfrak q\) of the original affine finite-type algebra, the integer \(\dim_xX\) is the minimum dimension of its open neighborhoods. Basic opens form a basis: if a neighborhood attains the minimum, some \(D(g)\) with \(g\notin\mathfrak q\) lies inside it, and its dimension is squeezed between that minimum and the chosen neighborhood dimension. Thus \(\dim S_g=\dim_xX\). The localization remains finite type, with presentation \(S_g=S[t]/(gt-1)\), retaining \(g\) and its inverse. The preceding induction applies to this actual ring. The zero ring has no such point, so no point case is missing.

\subsubsection{The refined map, every unit and a finite free strengthening}\label{algebra-proof-026-the-refined-map-every-unit-and-a-finite-free-strengthening}

Source 28073--28123. Retain the original prime \(\mathfrak q\subset A=k[x_1,\ldots,x_n]\) and \(r=\operatorname{trdeg}_k\kappa(\mathfrak q)\). If \(\mathfrak q=0\), then \(r=n\); take the identity map. Conversely a nonzero polynomial relation among the images of the \(n\) generators makes their field transcendence degree at most \(n-1\): a maximal algebraically independent subfamily has fewer than \(n\) elements, and every omitted generator is algebraic over its field. Thus \(r=n\) forces \(\mathfrak q=0\). For a nonzero prime choose the original \(0\ne f\in\mathfrak q\).

Apply the actual \(\sigma\) above and put \(h=\sigma(f)\), \(\mathfrak q^\sigma=\sigma(\mathfrak q)\). Retain \(h=a x_n^d+\sum_{j<d}c_j(x_1,\ldots,x_{n-1})x_n^j\), \(a\in k^\times\). The coefficient ring is \(k[x_1,\ldots,x_{n-1}]\), as corrected in OCC-00550 and OCC-12273. The source's monic wording requires multiplication by \(a^{-1}\), in addition to the automorphism. Rather than silently replacing \(f\), track both choices: its monic version is \(h_0=a^{-1}h\in\mathfrak q^\sigma\), while the following proof works with the unscaled \(h\) and keeps \(a\) in every relation. Changing the last source generator by \(Y_n\mapsto aY_n\) identifies the two finite maps, and \((Y_{r+1},\ldots,Y_n)\) is unchanged by this unit.

Define \(B=k[Y_1,\ldots,Y_n]\to A\) by \(Y_i\mapsto x_i\) for \(i<n\) and \(Y_n\mapsto h\). There is the exact presentation \[
 A\simeq B[T]\big/
 \left(T^d+\sum_{j<d}a^{-1}c_j(Y_1,\ldots,Y_{n-1})T^j-a^{-1}Y_n\right),
\] sending \(T\mapsto x_n\). Its inverse sends \(x_i\) to \(Y_i\), \(i<n\), and \(x_n\) to \(T\); the defining equation identifies the image of \(h\) with \(Y_n\), so both composites fix all generators. Monic division proves that \(A\) is free over \(B\) on \(1,T,\ldots,T^{d-1}\). The coefficient map is injective because a free module basis containing \(1\) makes its coefficient zero whenever that coefficient times \(1\) is zero.

The contraction \(Q=(B\to A)^{-1}(\mathfrak q^\sigma)\) contains \(Y_n\), so \(Q=\mathfrak pB+(Y_n)\) with \(\mathfrak p\subset k[Y_1,\ldots,Y_{n-1}]\) prime. The finite injection \(B/Q\to A/\mathfrak q^\sigma\) gives a finite extension of fraction fields, hence \(\operatorname{trdeg}_k\kappa(\mathfrak p)=r\). Inductively choose the refined finite map on \(n-1\) variables carrying \(\mathfrak p\) back to \((Z_{r+1},\ldots,Z_{n-1})\). Extend it polynomially by \(Z_n\mapsto Y_n\). The preimage of \(Q\) is now \((Z_{r+1},\ldots,Z_n)\): modulo \(Z_n,Y_n\) this is exactly the induction statement, and the last generator lies in the contracted ideal. Composing with \(B\to A\) and then \(\sigma^{-1}:A\to A\) carries this ideal to the original \(\mathfrak q\), not to a substituted prime.

This proves a stronger assertion: the refined finite map may be chosen injective and finite free of positive rank. The inductive map has a finite free basis by the stronger induction; polynomial extension preserves that basis. If it is \(b_1,\ldots,b_D\), the composite before \(\sigma^{-1}\) has basis \(b_jT^\ell\), \(1\leq j\leq D,\ 0\leq\ell<d\). Transport that exact basis by \(\sigma^{-1}\) for the original target. This proves finite freeness and injectivity at every step, including rank one at \(n=0\) or \(\mathfrak q=0\).

Write the final original images as \(a_i=\varphi(Z_i)\). Then \(a_{r+1},\ldots,a_n\) is a regular sequence in \(A\): the coordinate sequence in the polynomial source is regular, and freeness preserves each injection and each quotient. Its final quotient is nonzero since the original \(\mathfrak q\) contains the sequence. Moreover \(A/(a_{r+1},\ldots,a_n)\) is finite free of the same rank over \(k[Z_1,\ldots,Z_r]\). The prime \(\mathfrak q/(a_{r+1},\ldots,a_n)\) is minimal in this quotient: its contraction is zero, any smaller prime also contracts to zero, and incomparability for the integral extension prohibits strict inclusion. Thus the original prime is a minimal prime of this explicitly constructed regular-sequence quotient. No assertion that \(\mathfrak q\) itself is generated by that sequence has been substituted.

\subsubsection{Generic comparison over a domain without losing torsion}\label{algebra-proof-026-generic-comparison-over-a-domain-without-losing-torsion}

Source 28125--28161. Let the original injection \(R\to S\) be finite type with \(R\) a domain, \(K=\operatorname{Frac}R\), and fixed generators \(x_i\). The algebra \(S_K\) is nonzero, since \(R\to S\) is injective. The field construction supplies \(y_1,\ldots,y_d\) which are integer polynomials in those \(x_i\), so they are actual elements of \(S\). Any polynomial relation over \(R\) among these elements maps to a relation over \(K\); injectivity of the polynomial map over \(K\) and \(R\to K\) show all its original coefficients vanish. Thus \(R[y]\to S\) is injective. The source's change of dummy index from \(i\) to \(j\) refers to the same already specified finite family; OCC-00551 is an optional stylistic change.

Choose the source monic polynomials \(P_i(T)=T^{m_i}+\sum_{j<m_i}c_{ij}(y)T^j\in K[y][T]\) with \(P_i(x_i)=0\) in \(S_K\), necessarily \(m_i\geq1\). Choose nonzero \(g\in R\) clearing all coefficients \(c_{ij}\). The element \(gP_i(x_i)\) is then defined in \(S\) and maps to zero in \(S_K\); localization gives nonzero \(t_i\in R\) with \(t_i gP_i(x_i)=0\) in \(S\). Put \(f=g\prod_i t_i\), nonzero because \(R\) is a domain. Then every \(fP_i\) has coefficients in \(R[y]\), and \(fP_i(x_i)=0\) in \(S\).

Retain \(x_i'=fx_i\) and the exact polynomial \[
 Q_i(T)=f^{m_i}P_i(T/f)
       =T^{m_i}+\sum_{j<m_i} f^{m_i-j}c_{ij}(y)T^j .
\] Every coefficient is in \(R[y]\), since \(m_i-j\geq1\) and \(fc_{ij}\) already is. Evaluation in \(S\) gives \(Q_i(x_i')=f^{m_i-1}(fP_i(x_i))=0\). These formulas retain every factor, including those killing localization torsion. For \(S'=R[y,x_1',\ldots,x_n']\subset S\), repeated use of the monic equations shows that the products \(\prod_i(x_i')^{b_i}\), \(0\leq b_i<m_i\), generate \(S'\) as an \(R[y]\)-module. They are not claimed linearly independent.

The localized inclusion \(S'_f\to S_f\) is injective: if \(s'/f^a\) maps to zero, some \(f^b s'=0\) in \(S\), and the same equation holds in the actual subring \(S'\). It is onto because each original generator \(x_i=x_i'/f\) lies in its image. This gives both directions of the claimed identification, with every original denominator retained.

Injectivity \(R\to S\) does not make \(S\to S_K\) injective. For example \(R=k[t]\), \(S=R[z]/(tz)\) has injective coefficient map (split by \(z\mapsto0\)), but \(z\ne0\) in \(S\) and \(z=0\) in \(S_K\). Here the generic relation \(P(T)=T\) must be multiplied by \(t\); with \(f=t\), \(x'=tz=0\), \(S'=R\), and \(S'_t=S_t\). This verifies exactly why the source's annihilator step cannot be omitted.

\subsubsection{Dimension, transcendence degree and the zero-length chains}\label{algebra-proof-026-dimension-transcendence-degree-and-the-zero-length-chains}

Source 28177--28287. For a finite-type domain \(S\), take the source finite injection \(k[y_1,\ldots,y_d]\to S\), \(d=\dim S\). Tensor with the fraction field \(k(y)\). The resulting finite-dimensional domain over \(k(y)\) is a field, since multiplication by a nonzero element is an injective endomorphism of a finite-dimensional vector space. It contains \(S\) and the inverses of every nonzero element of \(S\), so it is its fraction field \(K\). This proves the required field extension is finite and gives \(\operatorname{trdeg}_kK=d\). The earlier exact prime-interval calculation gives the same dimension at every maximal ideal. The k-algebra wording at 28180 is corrected minimally.

If \(\mathfrak q\subsetneq\mathfrak q'\), work in the finite-type domain \(S/\mathfrak q\). A maximal chain in its proper prime quotient \(S/\mathfrak q'\) extends at its bottom by \(0\subsetneq\mathfrak q'/\mathfrak q\). Thus the dimension strictly decreases, and the just-proved equality with residue-field transcendence degree gives the source specialization inequality. Finiteness of these dimensions is essential to the strict integer comparison.

For a general finite-type \(S\), retain \(\mathfrak p\), the minimal primes \(\mathfrak q_j\subset\mathfrak p\), and \(r=\operatorname{trdeg}_k\kappa(\mathfrak p)=\dim S/\mathfrak p\). For every \(j\), the earlier polynomial-presentation chain formula gives \[
 \dim(S/\mathfrak q_j)
   =\dim (S/\mathfrak q_j)_{\mathfrak p/\mathfrak q_j}+r.
\] Taking the maximum over these actual components yields \[
 \dim_x\operatorname{Spec}S
       =\dim S_{\mathfrak p}+\operatorname{trdeg}_k\kappa(\mathfrak p).
\] This is precisely the source's concatenated-chain argument. If \(r=0\), the upper chain consists of \(\mathfrak p\) alone. If its lower interval has length \(s(j)=0\), then \(\mathfrak q_j=\mathfrak p\), and no \(\mathfrak p'_0\) is introduced. These are separate proof clarifications of the displayed chains.

For a surjection of finite-type \(k\)-algebras \(S'\to S\), the corresponding residue fields \(\kappa(\mathfrak p')\to\kappa(\mathfrak p)\) are isomorphic: localization and quotient give the same fraction field of \(S'/\mathfrak p'=S/\mathfrak p\). Subtract the two preceding formulas to get the exact stated difference of point dimensions and prime heights. The two other occurrences of k algebra(s), at 28224 and 28272, receive only the missing noun hyphens.

\subsubsection{Field extension and the common original fibre ring}\label{algebra-proof-026-field-extension-and-the-common-original-fibre-ring}

Source 28289--28387. If \(S=0\), then \(K\otimes_k S=0\) and both dimensions are \(-\infty\); this case needs no polynomial ring on \(d=-\infty\) variables. For \(S\ne0\), tensor the finite injection furnished above with \(K\). Flatness of the field extension preserves injectivity and finite generation, giving the original finite injection \(K[y_1,\ldots,y_d]\to K\otimes_kS\), and integral dimension equality proves the asserted global equality.

For the pointwise assertion, keep the source presentation \(A=k[x_1,\ldots,x_n]\to S=A/I\), \(B=K[x_1,\ldots,x_n]\), and \(S_K=B/IB\). Retain the original primes \(\mathfrak q'\subset A\), \(\mathfrak q_K'\subset B\) and their quotient primes \(\mathfrak q,\mathfrak q_K\). Both vertical localized maps are flat. Their common local fibre is exactly \[
 B_{\mathfrak q_K'}/\mathfrak q'B_{\mathfrak q_K'}
 \ \simeq\
 (S_K)_{\mathfrak q_K}/\mathfrak q(S_K)_{\mathfrak q_K}.
\] The map sends \([b/u]\) to \([\bar b/\bar u]\); its inverse lifts quotient representatives. It is well defined and bijective because \(I\subset\mathfrak q'\), so quotienting by \(IB\) before \(\mathfrak q'B\) has the same final kernel. All denominators avoid the specified primes.

Writing the fibre dimension as \(h\), the local base-plus-fibre equality gives \[
 \operatorname{ht}(\mathfrak q_K')
     =\operatorname{ht}(\mathfrak q')+h,\qquad
 \operatorname{ht}(\mathfrak q_K)
     =\operatorname{ht}(\mathfrak q)+h.
\] The source's quotient codimension formula subtracts these two equalities. The point dimension of either ambient affine \(n\)-space is \(n\), including at nonclosed points, by its single component and the earlier point-dimension result. Hence \(\dim_xX=\dim_{x_K}X_K\). This proves the source cancellation through the exact same fibre ring, not merely through equally named dimensions. The duplicated ``Denote'' reports receive ``Denote by'' at 28343.

Combining the local dimension equality and the pointwise formula gives \[
 h=\dim(S_K)_{\mathfrak q_K}-\dim S_{\mathfrak q}
   =\operatorname{trdeg}_k\kappa(\mathfrak q)
      -\operatorname{trdeg}_K\kappa(\mathfrak q_K).
\] The missing finite verb at 28377 is corrected as reported twice. A prime minimal over \(\mathfrak qS_K\) exists because \(S_K\) is Noetherian and that ideal is proper. Properness follows from faithful flatness of the field base change. Such a prime contracts to \(\mathfrak q\): going down supplies a smaller prime over \(\mathfrak q\), and it still contains \(\mathfrak qS_K\), so minimality forces equality. The localized fibre then has only its minimal maximal prime and dimension zero. This proves the final source existence assertion with the required contraction checked.

\subsubsection{Exact maximum fibre dimension and both height-jump extremes}\label{algebra-proof-026-exact-maximum-fibre-dimension-and-both-height-jump-extremes}

Let \(F/k\) be the original finitely generated residue field \(F=\kappa(\mathfrak q)\), put \(r=\operatorname{trdeg}_kF<\infty\), and let \(K/k\) be arbitrary, of transcendence degree \(s\), possibly infinite. Then the actual tensor ring, including its possible nilpotents, satisfies \[
 \dim(K\otimes_kF)=\min(r,s).
\] Here is a direct proof. Choose an actual transcendence basis \(t_1,\ldots,t_r\) in \(F\). The extension \(F/k(t_1,\ldots,t_r)\) is finite: \(F\) is finitely generated as a field, and the remaining generators are algebraic. If \(D\) is its degree, tensoring its vector space basis gives a finite free, hence integral and injective, comparison \[
 A=(k[t_1,\ldots,t_r]\setminus\{0\})^{-1}K[t_1,\ldots,t_r]
      \ \longrightarrow\ K\otimes_kF .
\] It is obtained from \((K\otimes_k k(t))\otimes_{k(t)}F\simeq K\otimes_kF\); \((a\otimes b)\otimes c\mapsto a\otimes bc\) and \(a\otimes c\mapsto(a\otimes1)\otimes c\) are inverse. No reducedness or separability assumption is introduced. Integral dimension equality reduces only this dimension calculation to \(A\); the original tensor ring remains present.

Every prime \(P\subset K[t_1,\ldots,t_r]\) surviving in \(A\) contracts to zero in \(k[t_1,\ldots,t_r]\), so its residue field \(L\) contains \(k(t_1,\ldots,t_r)\). The earlier affine-domain dimension formula gives \(\operatorname{ht}P=r-\operatorname{trdeg}_K L\). This is at most \(r\). If \(s<\infty\), the actual field tower gives \(\operatorname{trdeg}_k L=s+\operatorname{trdeg}_K L\geq r\), so \(\operatorname{ht}P\leq s\). The entire interval below \(P\) survives the localization, giving the upper bound \(\dim A\leq\min(r,s)\).

For the lower bound put \(m=\min(r,s)\) and choose algebraically independent \(u_1,\ldots,u_m\in K\) over \(k\). In \(K[t_1,\ldots,t_r]\), take \(P=(t_1-u_1,\ldots,t_m-u_m)\). Its quotient is \(K[t_{m+1},\ldots,t_r]\). Evaluation of \(k[t_1,\ldots,t_r]\) into this quotient is injective: expand in the last \(r-m\) variables and use algebraic independence of the \(u_i\) on each coefficient. Thus \(P\) is disjoint from the localization set. Its height is \(m\), either by the exact dimension formula for this polynomial quotient or by its successive coordinate prime chain and its \(m\)-generator local maximal ideal bound. All those lower primes survive. This proves the lower bound, and hence the displayed equality. For \(r=0\) or \(m=0\), the argument uses empty transcendence bases and the zero prime directly.

Consequently every source height jump \(h\) obeys \(0\leq h\leq\min(r,s)\), and both extremes occur among primes over the fixed \(\mathfrak q\). For zero, use a minimal prime of the nonzero Noetherian tensor ring. For the maximum, a finite-dimensional spectrum attains its integer dimension: a finite strict chain of maximal possible length has a last prime whose local dimension equals that length. The prime correspondence between \(\operatorname{Spec}(K\otimes_kF)\) and primes of \(S_K\) contracting to \(\mathfrak q\) preserves their local fibre rings. Thus both choices really lie over the original \(\mathfrak q\). For an algebraic extension \(K/k\), \(s=0\), so every height jump is zero.

As an exact example retaining distinct base and point dimensions, take \(S=k[x]\), \(\mathfrak q=0\), \(K=k(u)\). Both primes \(0\) and \((x-u)\) of \(K[x]\) contract to zero in \(k[x]\); their local heights are respectively zero and one, while the topological dimension at both points remains one. This realizes both extremes without suggesting that point dimension and local height are identical.

\subsubsection{The graded vertex, its Hilbert function and its full completion}\label{algebra-proof-026-the-graded-vertex-its-hilbert-function-and-its-full-completion}

Source 28402--28434. Keep the original finitely generated standard graded \(k\)-algebra \(S\), \(S_0=k\), and \(\mathfrak m=S_+\). Its quotient is \(k\), so \(\mathfrak m\) is maximal. Every minimal prime is homogeneous by the source homogeneous-prime lemma. A proper homogeneous ideal contains no nonzero degree-zero scalar, so every such minimal prime lies in \(\mathfrak m\). All components therefore pass through the vertex. The earlier point dimension formula gives \(\dim S=\dim_{\mathfrak m}\operatorname{Spec}S=\dim S_{\mathfrak m}\).

Standard generation gives the exact ideals \[
 \mathfrak m^d=\bigoplus_{j\geq d}S_j.
\] Every product of \(d\) positive-degree elements has degree at least \(d\). Conversely each degree-\(j\) monomial in the degree-one generators, \(j\geq d\), factors as a product of \(d\) positive-degree monomials; linear combinations give the reverse inclusion. Thus the degree-\(d\) projection identifies \(\mathfrak m^d/\mathfrak m^{d+1}\) with \(S_d\), including \(d=0\).

Put \(R=S_{\mathfrak m}\), \(\mathfrak n=\mathfrak mR\). For \(u\notin\mathfrak m\), write its original homogeneous decomposition \(u=c+v\), \(c\in k^\times\), \(v\in\mathfrak m\). On \(\mathfrak m^d/\mathfrak m^{d+1}\), multiplication by \(u\) is multiplication by \(c\), with inverse \(c^{-1}\). This proves that localization gives an isomorphism onto \(\mathfrak n^d/\mathfrak n^{d+1}\), with every denominator accounted for. The source local Hilbert function is exactly \(\operatorname{length}_R(\mathfrak n^d/\mathfrak n^{d+1})\), as read at 13994--14015; it is not the cumulative length. Since these quotients are vector spaces over \(k\), that length is \(\dim_k S_d\). This proves the claimed equality of functions.

If the eventual polynomial \(P\) is nonzero, its leading coefficient is positive (its values are eventual nonnegative dimensions), and summing from zero through \(n\) increases its degree by one. The source local dimension theorem applied to that cumulative polynomial therefore gives \(\dim R=\deg P+1\). If \(P=0\), all sufficiently high \(S_d\) vanish, so \(S\) is finite dimensional, \(\mathfrak m\) is nilpotent and \(R\) is nonzero Artinian of dimension zero. This gives \(\deg(0)+1=-1+1=0\) exactly. No subtraction of infinities or negative binomial indices is used. The definitions and the prior group 384's zero-degree check agree.

In fact the preceding degree maps are multiplicative and prove the complete graded algebra identification \[
 S\longrightarrow\operatorname{gr}_{\mathfrak n}R,\qquad
 a\in S_d\longmapsto a/1\bmod\mathfrak n^{d+1}.
\] Every degree map is bijective by the explicit localization inverse above, and multiplication preserves degrees, so this is an isomorphism of graded \(k\)-algebras.

For every \(r\geq1\), the full localized quotient map \(S/\mathfrak m^r\to R/\mathfrak n^r\) is also an isomorphism. Indeed the inverse of \(u=c+v\) modulo \(\mathfrak m^r\) is the exact finite expression \[
 c^{-1}\sum_{j=0}^{r-1}(-c^{-1}v)^j.
\] Multiplication by \(c+v=c(1+c^{-1}v)\) gives \(1-(-c^{-1}v)^r\), which equals one in that quotient. This supplies an inverse for every localized denominator without assuming one in the original ring. The quotient is the direct sum of its original degrees \(0,\ldots,r-1\). Passing to compatible quotient systems therefore gives inverse maps \[
 \widehat{R}\ \simeq\ \prod_{d\geq0}S_d .
\] On the right multiplication is the actual convolution \((a b)_d=\sum_{i=0}^d a_i b_{d-i}\); this sum is finite in each degree. A sequence of coefficients maps to its finite truncations, and a compatible system maps to its uniquely stabilized coefficients in each degree. These operations are inverse, preserve addition and convolution, and identify the respective adic topologies. All homogeneous nilpotents and their products are retained. This is a proved comparison with the original completion, not a replacement of the graded ring or its Hilbert function.

\subsubsection{Corpus and continuation boundaries}\label{algebra-proof-026-corpus-and-continuation-boundaries}

The two bounded corpus queries return eight routing matches across both layers and five distinct canonical IDs. They include physics records, local programme notes and a PDF-primary algebra volume. None is read or adopted as a theorem for this batch; no PDF fallback is used. The existing primary Stacks source and earlier source-bound proofs supply the exact calculations above. No new claim of human-literature novelty or exhaustive topic coverage is made. The previous original Heinrich and other bounded reading records are retained unchanged in the topic route.

All seventeen reports through 28447 receive dispositions. Required edits are minimal; the two stylistic reports remain unapplied. Continue at 28448, Generic flatness. Lookahead through 28460 has been read but remains unadjudicated.
